\documentclass[12pt]{article}
\usepackage[utf8]{inputenc}
\usepackage[T1]{fontenc}
\usepackage{hyperref}
\title{Tonkawa Language Resources for K-12 Curriculum Development}
\author{Compiled by Claude Sonnet 5 (claude-sonnet-5) for Andr\'as Kornai}
\date{}
\begin{document}
\maketitle

\section{Introduction}
Tonkawa (ISO 639-3 tqw; Glottocode tonk1249) is a dormant language isolate formerly spoken in Texas and Oklahoma. Wikipedia reports no first-language speakers since the mid-20th century, and the publisher's page for the Hoijer texts says it has been extinct for more than 75 years. Wikipedia also mentions revitalization efforts, but this was not independently confirmed. The Tonkawa Tribe of Oklahoma has produced language books.

\section{Primary Educational Resources}
\begin{itemize}
\item Tonkawa tribal language books from the Tonkawa Tribe Language Restoration and Presentation Program Series (2015), including an elementary outline of Tonkawa grammar by former chief Don Patterson. Ten copies were donated to the Northern Oklahoma College Cultural Engagement Center. See: \url{https://www.noc.edu/noc-news/cultural-engagement-center-receives-tonkawa-language-cultural-books/}
\item Local news coverage of the same donation: \url{https://www.blackwelljournaltribune.net/articles/11713/view/tonkawa-tribe-donates-language-books-to-noc}
\end{itemize}

\section{Linguistic and Reference Materials}
\begin{itemize}
\item Hoijer's Tonkawa texts, edited and translated by Thomas R. Wier (University of Oklahoma Press), with a grammatical description: \url{https://www.oupress.com/9780806158990/}
\item Tonkawa texts: a new linguistic edition (library record): \url{https://search.ub.tum.de/vufind/Record/DE-604.BV044854005?lng=en}
\item Tonkawa language (Wikipedia): \url{https://en.wikipedia.org/wiki/Tonkawa_language}
\item Tonkawa numerals (Max Planck Institute for Evolutionary Anthropology): \url{https://lingweb.eva.mpg.de/channumerals/Tonkawa.htm}
\end{itemize}

\section{Religious and Community Resources}
None found beyond the tribal language program above.

\section{Digital and Audio Resources}
No audio or digital teaching resources were found. No dictionary was found in the search results.

\section{Curriculum Development Context}
A heritage curriculum would rest on the Hoijer texts (the main surviving oral literature) and on the tribe's own language books. Collaboration with the Tonkawa Tribe of Oklahoma would be essential. No school programme was found.
\end{document}